\chapter{Fazit und Ausblick}
Die vorliegende Arbeit zeigt, dass die Entwicklung und Implementierung eines KI-basierten Live-Meetingassistenten auf Basis moderner Speech-to-Text-Systeme und Large Language Models technisch realisierbar ist und ein Potenzial zur Unterstützung wissensintensiver Arbeitsprozesse besitzt. Der entwickelte Prototyp zeigt, wie gesprochene Inhalte während laufender Meetings automatisiert erfasst, analysiert und aufbereitet werden können. Dadurch kann, in der Theorie, die Nachbereitung von Besprechungen erleichtert und eine schnellere Orientierung über zentrale Themen, Aufgaben und Entscheidungen ermöglicht.

Ein zentrales Ergebnis dieser Arbeit ist die Konzeption und Entwicklung einer modularen und skalierbaren Systemarchitektur, die eine inkrementelle Verarbeitung kontinuierlicher Transkriptdaten ermöglicht. Durch die klare Trennung von Transkription, Backend-Verarbeitung, LLM-Analyse und Frontend-Darstellung konnte eine robuste Pipeline geschaffen werden, die sich flexibel erweitern und in die bestehende Unternehmensinfrastrukturen von Notisent integrieren lässt. Der Einsatz etablierter Technologien wie Django sowie die Anbindung externer KI-Dienste ermöglichen eine wartbare und zukunftsfähige Implementierung. Gleichzeitig erlaubt die Verwendung offener Modellarchitekturen eine Anpassung an zukünftige Entwicklungen.

In den durchgeführten Testmeetings wurden Zusammenfassungen typischerweise innerhalb von Sekunden nach Eingang des Transkriptsegments erzeugt. Die Ergebnisse zeigen jedoch ebenfalls, dass die Qualität der inkrementellen Zusammenfassungen noch nicht das Niveau einer nachträglichen Gesamtauswertung erreicht. Insbesondere Kontextverluste, Transkriptionsfehler sowie die fortlaufende Kompression im Rahmen der Rolling-Summary-Logik führen zu Informationsverlusten und inkonsistenten Gewichtungen. Der Assistent in der aktuellen Form stellt somit höchstens eine Unterstützung dar, ersetzt jedoch noch keine manuelle Dokumentation.

Für einen produktiven Einsatz bestehen daher weiterhin Herausforderungen, welche vorher gelöst werden müssen. Besonders die Robustheit gegenüber fehlerhaften Transkripten, eine stabilere Kontextverwaltung über längere Gesprächsdauern sowie eine zuverlässige Extraktion von Aufgaben, Verantwortlichkeiten und Entscheidungen erfordern weitere Forschungs- und Entwicklungsarbeit. Neben der Optimierung der bestehenden Zusammenfassungslogik, wie in Kapitel \ref{optimierung} beschrieben, bieten sich zusätzliche Vorverarbeitungsschritte wie semantische Segmentierung, strukturierte Informationsextraktion oder klassische NLP-Verfahren zur Qualitätssteigerung an.

Zukünftige Arbeiten können darüber hinaus die Integration weiterer Datenquellen und Funktionen untersuchen. Dazu zählen beispielsweise die Verknüpfung mit Kalendern, Aufgabenmanagement-Systemen oder Dokumentenablagen, um Kontextinformationen automatisch einzubeziehen. Auch der Einsatz eigener Speech-to-Text-Modelle zur Verbesserung der Datenhoheit sowie die Erweiterung um multimodale Daten wie Audio- oder Bildschirmaufzeichnungen stellen vielversprechende Ansätze dar, welche auch von Interesse von Notisent sind. Schließlich ist die systematische Evaluation der Nutzerakzeptanz und des tatsächlichen Produktivitätsgewinns im Arbeitsalltag ein wichtiger nächster Schritt.

Insgesamt zeigt die Arbeit, dass ein LLM-gestützter Live-Meetingassistent heute technisch umsetzbar ist und bereits einen Nutzen bietet. Gleichzeitig wird deutlich, dass weitere Optimierungen notwendig sind, um eine zuverlässige, vollständig automatisierte Meetingdokumentation zu erreichen. Der entwickelte Prototyp bildet hierfür eine   Grundlage und schafft die Voraussetzungen für zukünftige Forschung und praxisnahe Weiterentwicklungen.


\chapter{Danksagung und Hilfsmittel}
\label{sec:hilfsmittel}
An dieser Stelle möchte ich mich bei allen bedanken, die mich während meiner Masterarbeit zum Thema \textit{Entwicklung eines KI-basierten Meetingassistenten zur Darstellung und Aufbereitung gesprochener Inhalte} unterstützt und begleitet haben.

Mein besonderer Dank gilt meinen Betreuern Prof. Dr. Roman Grothausmann und Henri Reumschüssel für ihre fachliche Expertise, wertvollen Hinweise und die kontinuierliche Unterstützung während des gesamten Forschungsprozesses. Ohne ihre konstruktive Kritik und ihr Engagement wäre diese Arbeit nicht in der vorliegenden Form möglich gewesen.

Ein weiterer Dank gilt dem Team von Notisent für den wertvollen Austausch, die gemeinsame Zeit im Büro, die mich bei meiner Arbeit unterstützt hat, sowie für die Möglichkeit, an den Meetings zur Datenerhebung teilzunehmen. Darüber hinaus möchte ich mich bei Christian Teich für die Teilnahme am Interview zu dem Produkt bedanken. Nicht zuletzt möchte ich mich bei noch Sürig Ohainski, Sonja Fiehn und Klara Jakubzik fürs Korrekturlesen bedanken.

Als Hilfsmittel wurden vor allem die Chatbot-Software \textit{ChatGPT} und der Code-Assistant \textit{Github Copilot}, welcher als Plugin in Visual Studio Code verwendet wurde, eingesetzt. \textit{ChatGPT} wurde überwiegend zur Ideenfindung als \textit{Sparrings-Partner}, zur Recherche und zum Überarbeiten von Textabschnitten verwendet. Die Tools \textit{Github Copilot} und \textit{ChatGPT} wurden beide als KI-Assistenten zum Programmieren verwendet. \textit{ChatGPT} wurde in der Version 5.2 genutzt, \textit{Github Copilot} mit den Modellen \textit{Claude Sonnet 4.5}, \textit{GPT-5.0} sowie \textit{GPT-4.1} verwendet.
Zur Erstellung von Grafiken und Abbildungen wurde die Online-Software \textit{Draw.io} verwendet.